% !PRINTED 575
\weberpage{575}

%<<CONT
satisfy. This equation has only real roots if $f$ is even, because then $\tfrac{1}{2}(n - 1)$ is a multiple of $e$, and consequently $r$ and $r^{-1}$ occur in the same period. But if $f$ is odd, all the roots are imaginary.

It is now of the highest interest to investigate this equation of degree $e$ for individual particular values of $e$, without any special assumption about the prime $n$ beyond the requirement that $n - 1$ shall be divisible by $e$.

We consider the first special cases, and assume to begin with $e = 2$, which is admissible for every odd prime $n$; here $\eta, \eta_1$ are thus the two periods of $\tfrac{1}{2}(n - 1)$ terms, which we now wish to denote by $A, B$, so that
\[
\begin{aligned}
A &= r + r_2 + r_4 + \cdots + r_{n-3}\\
B &= r_1 + r_3 + r_5 + \cdots + r_{n-2}.
\end{aligned}
\]
The $r, r_2, r_4 \ldots$ have the exponents $g^0, g^2 \ldots g^{n-3}$, i.\,e. the exponents of $g$ are even numbers. The exponents of $r$ are therefore the \gesperrt{quadratic residues} of $n$ (\S.\,138). Likewise the exponents of $r$ in the sum $B$ are the non-residues. If, then, we denote the residues by $a$ and the non-residues by $b$, we have
\weq{1}{A = \Sigma r^a, \quad B = \Sigma r^b,}
and $(-1, r) = A - B$. These expressions $A, B$, are called the \gesperrt{Gauss sums}.

If in the formula \S.\,169, (10) we make the assumption
\[
e = 2, \quad \mu = f = \frac{n-1}{2}, \quad \lambda = 1, \quad \alpha = -1,
\]
there results
\weq{2}{A - B = \pm \sqrt{(-1)^{\frac{n-1}{2}}\, n},}
while on the other hand
\[
A + B = -1,
\]
hence
\weq{3}{2A = -1 \pm \sqrt{(-1)^{\frac{n-1}{2}}\, n}, \quad 2B = -1 \mp \sqrt{(-1)^{\frac{n-1}{2}}\, n}.}
The sign that is to be given to the root depends on the choice of $r$. But once $r$ has been fixed upon, the sign is completely determined. Its ascertainment presents peculiar difficulties, which gave occasion to a separate memoir by \gesperrt{Gauss}\textsuperscript{1}).

\noindent\rule{3cm}{0.4pt}

\noindent\textsuperscript{1}) Summatio quarundam serierum singularium. \gesperrt{Gauss} Werke, Bd.\ II. In another way, by application of the higher analysis, the sign
%CONT>>
